% =====================================================================
%  第 5 节  结语（原书 p.25）
% =====================================================================
\section{结语}\label{sec:5}

\begin{takeaway}
\begin{itemize}
  \item \textbf{物理图像}：第~\ref{sec:2.4} 节把余维一闭链纳入态射，相当于给对象加上一个平坦联络（局部系统）；把对象换成 $\nu_a$ 个副本（$T^*X$ 语言中的“堆叠”），就得到系数在 $GL(\nu_a)$ 中的非阿贝尔平坦联络。
  \item \textbf{数学工具}：Fukaya 范畴的“堆叠”语言、分块矩阵代数、Morse--Bott 函数与 Morse--Bott 链。
  \item \textbf{主要结论}：全文两个证明都没有用到全纯圆盘这一概念；构造还可推广到 $F_a-F_b$ 仅为 Morse--Bott 的情形，此时 $\Mor(F_a,F_b)$ 取为 $\mathrm{Crit}(F_a-F_b)$ 上的 Morse--Bott 链。
\end{itemize}
\end{takeaway}

按第~\ref{sec:2.4} 节的方式把余维一的闭链纳入进来，可以解释为给一个对象加上平坦联络（局部系统）。其实还可以考虑把对象 $F_a$ 换成 $F_a$ 的 $\nu_a$ 个副本——用 $T^*X$ 的 Fukaya 范畴的语言说，就是 $\nu_a$ 个 Lagrange 子流形 $L_a$ 的“堆叠”。用第~\ref{sec:3} 节的语言，这相当于把 $N\times N$ 矩阵构成的系数代数，换成以 $\nu_a\times\nu_b$ 为块大小的分块矩阵代数。特别地，由此可以对一个对象 $F_a$ 配上 $X$ 上（或用 $T^*X$ 的语言，$L_a$ 上）系数在 $GL(\nu_a)$ 中的非阿贝尔平坦联络。

% [Tutor Note: 余维一闭链与平坦联络的对应来自式 (5)：$m(Z_1,Z_2)=Z_1\cap Z_2$ 是链的带符号相交。取 $Z$ 为余维一闭链时，相交数给出绕该闭链一周的相位，这正是和乐；把系数换成 $GL(\nu_a)$ 矩阵，和乐就成为非阿贝尔的。作者没有展开这个对应，正文只用了「可以解释为」的说法。]

% [Tutor Note: “堆叠”是 $T^*X$ 的 Fukaya 范畴中的标准图像：一个 Lagrange 子流形取 $\nu$ 个重合副本，其上的规范场联络自然取 $GL(\nu)$ 值。注意这里的副本不是“扰动开来的”，而是重合的；若真的扰动开，全纯圆盘会重新出现，反而不属于本文的框架。]

注意，无论是在 $A_\infty$ 关系式的第一个证明（用有效作用量）中，还是在第二个证明（用拓扑量子力学）中，我们都没有用到全纯圆盘这一概念。

我们还要指出，本文的构造可以推广到如下情形：函数差 $F_a-F_b$ 只要求是 Morse--Bott 的，而不必是 Morse 的。此时临界集 $\mathrm{Crit}(F_a-F_b)$ 允许是不同维数的子流形之并，而不限于孤立点集。于是把 $\Mor(F_a,F_b)$ 取为 $\mathrm{Crit}(F_a-F_b)$ 上的奇异链，并配以 Morse--Bott 微分；该微分把链上的边界算子与临界流形之间梯度流的数据结合起来。注意，在这一处理方式下，无需单独讨论对象 $F_a$ 到自身的态射——它是上述一般构造的一个特例（事实上 $\mathrm{Crit}(F_a-F_a)=X$，因此此时的 Morse--Bott 链就是 $X$ 上的奇异链）。

% [Tutor Note: 末句把 §2.1 的“特设”模型解释成一般构造的推论：$F_a-F_a\equiv0$ 的临界集是整个 $X$，故自同态的 Morse--Bott 链复形就是 $X$ 上的奇异链，与式 (2) 一致。但要注意 Morse--Bott 微分比链的边界算子多一项——临界流形之间的梯度流；在 $a=b$ 的特例里这一项退化，所以 (2) 只需边界算子。]

\endinput
